\documentclass[a4paper,11pt]{article}
\usepackage[T1]{fontenc}
\usepackage[utf8]{inputenc}
\usepackage[italian]{babel}
\usepackage{amsmath,amssymb}
\usepackage{geometry}
\geometry{margin=2.5cm}

\title{PDDL: Domino}
\date{}

\begin{document}
\maketitle

\subsection*{1.1 Modellazione del Dominio} 

\paragraph{Predicati}
\[ P = \{\text{Tessera}/1, \text{Numero}/1, \text{HaValori}/3, \text{Disponibile}/1, \text{EstremitaAperta}/1, \text{TavoloVuoto}/0\} \]

\paragraph{Segnature delle costanti}
Definiamo le costanti (valori numerici disponibili):
\[ C = \{1, 2, 3, 4, 5, 6\} \]

\paragraph{Azioni}

\begin{itemize}
    \item \textbf{IniziaGioco}$(t, n_1, n_2)$
    \begin{itemize}
        \item $Pre: \text{Tessera}(t) \land \text{Numero}(n_1) \land \text{Numero}(n_2) \land \text{TavoloVuoto}() \land \text{Disponibile}(t) \land \text{HaValori}(t, n_1, n_2)$
        \item $Eff: \neg \text{TavoloVuoto}() \land \neg \text{Disponibile}(t) \land \text{EstremitaAperta}(n_2)$
    \end{itemize}
    \vspace{0.2cm}
    \item \textbf{IniziaGiocoRuotata}$(t, n_1, n_2)$
    \begin{itemize}
        \item $Pre: \text{Tessera}(t) \land \text{Numero}(n_1) \land \text{Numero}(n_2) \land \text{TavoloVuoto}() \land \text{Disponibile}(t) \land \text{HaValori}(t, n_1, n_2)$
        \item $Eff: \neg \text{TavoloVuoto}() \land \neg \text{Disponibile}(t) \land \text{EstremitaAperta}(n_1)$
    \end{itemize}
    \vspace{0.2cm}
    \item \textbf{GiocaTessera}$(t, n_{match}, n_{new})$
    \begin{itemize}
        \item $Pre: \text{Tessera}(t) \land \text{Numero}(n_{match}) \land \text{Numero}(n_{new}) \land \text{Disponibile}(t) \land \text{EstremitaAperta}(n_{match}) \land \text{HaValori}(t, n_{match}, n_{new})$
        \item $Eff: \neg \text{Disponibile}(t) \land \neg \text{EstremitaAperta}(n_{match}) \land \text{EstremitaAperta}(n_{new})$
    \end{itemize}
    \vspace{0.2cm}
    \item \textbf{GiocaTesseraRuotata}$(t, n_{match}, n_{new})$
    \begin{itemize}
        \item $Pre: \text{Tessera}(t) \land \text{Numero}(n_{match}) \land \text{Numero}(n_{new}) \land \text{Disponibile}(t) \land \text{EstremitaAperta}(n_{match}) \land \text{HaValori}(t, n_{new}, n_{match})$
        \item $Eff: \neg \text{Disponibile}(t) \land \neg \text{EstremitaAperta}(n_{match}) \land \text{EstremitaAperta}(n_{new})$
    \end{itemize}
\end{itemize}

\newpage

\subsection*{1.2 Istanziazione specifica}

\paragraph{Istanziazione (Oggetti)}
\[ F = \{ t_1, t_2, t_3, t_4, t_5, t_6 \} \]

\paragraph{Stato Iniziale}
\[
\begin{aligned}
I = & \text{Tessera}(t_1) \land \text{Tessera}(t_2) \land \text{Tessera}(t_3) \land \text{Tessera}(t_4) \land \text{Tessera}(t_5) \land \text{Tessera}(t_6) \land \\
& \text{Numero}(1) \land \text{Numero}(2) \land \text{Numero}(3) \land \text{Numero}(4) \land \text{Numero}(5) \land \text{Numero}(6) \land \\
& \text{HaValori}(t_1, 1, 2) \land \text{HaValori}(t_2, 1, 2) \land \text{HaValori}(t_3, 1, 3) \land \\
& \text{HaValori}(t_4, 3, 1) \land \text{HaValori}(t_5, 2, 4) \land \text{HaValori}(t_6, 2, 5) \land \\
& \text{Disponibile}(t_1) \land \text{Disponibile}(t_2) \land \text{Disponibile}(t_3) \land \\
& \text{Disponibile}(t_4) \land \text{Disponibile}(t_5) \land \text{Disponibile}(t_6) \land \\
& \text{TavoloVuoto}()
\end{aligned}
\]

\paragraph{Obiettivo}
\[
\begin{aligned}
G = & \neg \text{Disponibile}(t_1) \land \neg \text{Disponibile}(t_2) \land \neg \text{Disponibile}(t_3) \\
    & \land \neg \text{Disponibile}(t_4) \land \neg \text{Disponibile}(t_5) \land \neg \text{Disponibile}(t_6)
\end{aligned}
\]

\end{document}